% Condensed Appendix D; expanded execution/numerical records are in the artifact.
\subsection{Experimental Quantities and Aggregation}
\label{app:experimental-conventions}

Table~\ref{tab:d-experiment-overview} identifies the operational targets; Appendix~\ref{app:operational-matrices} defines the matrix constructions. Real-data audits use nuisance ridge $10^{-3}$ and stabilization $10^{-8}$; synthetic audits use $\varepsilon=10^{-12}$. Certificates use $\alpha=0.95$. Primary predefined-group coordinates use training moments; prospective normalization is specified below. Reported evaluation Loss-CVaR averages group-mean NLLs at or above their empirical 95th percentile, weighting groups equally; synthetic experiments use environment-mean NLLs. The Loss-CVaR training baseline averages the largest $\lceil0.1n\rceil$ individual-example NLLs.

\begin{table}[htbp]
\centering\small
\setlength{\tabcolsep}{3pt}
\caption{Diagnostic targets and replication units.}
\label{tab:d-experiment-overview}
\begin{tabular}{@{}p{.26\linewidth}p{.61\linewidth}c@{}}
\toprule
Experiment & Representation and reported quantity & Seeds \\
\midrule
Prospective Waterbirds & Fine-tuned MobileNetV3-small; background-probe surviving fraction & 22--26 \\
Synthetic & Normalized task/nuisance coordinates; inverse-information tail and certificate & 11--15 \\
Hidden Waterbirds slices & Frozen MobileNetV3-small and fixed CLIP direction; shrunken predictive information & 17--21 \\
CivilComments groups & TF-IDF logistic classifier; feature--Schur certificate & 13--17 \\
Waterbirds groups & Frozen MobileNetV3-small and logistic fitting; feature--Schur certificate & 17--21 \\
\bottomrule
\end{tabular}
\par\smallskip
\end{table}


\paragraph{Uncertainty.}
Prospective, synthetic, predefined-group, and frozen-feature training summaries report nominal 95\% normal-approximation half-widths $1.96s/\sqrt5$, with sample standard deviation $s$. Hidden-slice and group-free comparisons use five-seed Student-$t$ intervals; contrasts are paired within seed. The end-to-end frontier reports standard deviations, and the conditional Gaussian/Wald audit is separate (Appendix~\ref{app:implemented-audit}). Hidden-gate bootstrap intervals condition on the frozen gate and recompute evaluation weights (Appendix~\ref{app:hidden-slice-discovery}). Acquisition retains means without recovered per-seed intervals.

Seeds describe fitting/partition or control variation on fixed benchmark splits, not independent datasets. The frozen Waterbirds classifier and primary diagnostic are deterministic after feature extraction and logistic fitting; only its random controls and reserve anchor vary. Prospective intervals use training seeds: 23,176 composites share 5,794 birds, with two donors per bird and arm. Accuracy changes are percentage points; information retains its original units.

\subsection{Prospective Nuisance Intervention}
\label{app:prospective-intervention}

\paragraph{Training and frozen selection.}
ImageNet-initialized MobileNetV3-small~\citep{russakovsky2015imagenet,howard2019searching} is fully fine-tuned on Waterbirds~\citep{sagawa2019distributionally} for seeds 22--26: 50 epochs, batch size 32, learning rate $10^{-4}$, and bf16 on an H100. A validation-fitted temperature and class-balanced logistic background probe use the pooled representation. Five probe folds are stratified by validation group (3,000 maximum iterations); features use validation standardization. Validation nuisance scores are out of fold and use validation moments; test scores use the full-validation probe and test moments. Task logits use training moments. Test normalization never selects the slice. The rule is
\begin{equation}
\begin{aligned}
\mathcal G_{\mathrm{eligible}}=\{g:\;&\operatorname{Acc}_{v}(g)\ge\operatorname{Acc}_{v},\quad
\operatorname{NLL}_{v}(g)\le\operatorname{median}_{h}\operatorname{NLL}_{v}(h)\},\\
g^\star&\in\underset{g\in\mathcal G_{\mathrm{eligible}}}{\arg\min}
\lambda_{\min}(\widehat F_g^{\rid})/\widehat J_{TT,g}.
\end{aligned}
\label{eq:d-prospective-selection}
\end{equation}
Here $v$ denotes validation and eligibility uses uncalibrated NLL. Temperature scaling~\citep{guo2017calibration} minimizes full-validation binary NLL over 400 log-spaced temperatures in $[0.05,20]$, choosing the smaller temperature at ties. The checkpoint, temperature, probe, slice and audit constants are frozen before scoring; checkpoint mismatches abort evaluation. Table~\ref{tab:d-validation-performance} shows that LL (landbird/land) and WW (waterbird/water) qualify; WW is selected in all five seeds.

\begin{table}[htbp]
\centering\small
\setlength{\tabcolsep}{3pt}
\caption{Validation selection. Means $\pm$ nominal 95\% normal-approximation half-widths across seeds 22--26; accuracy is a fraction. Information uses uncentered feature moments.}
\label{tab:d-validation-performance}
\label{tab:d-validation-information}
\textit{(a) Performance and eligibility}\par\smallskip
\begin{tabular}{@{}lrccc@{}}
\toprule
Group & $n$ & Accuracy & NLL & Eligible \\
\midrule
Landbird/land & 467 & $0.990\pm0.005$ & $0.038\pm0.016$ & Yes \\
Landbird/water & 466 & $0.703\pm0.019$ & $1.670\pm0.156$ & No \\
Waterbird/land & 133 & $0.390\pm0.021$ & $7.718\pm0.800$ & No \\
Waterbird/water & 133 & $0.920\pm0.013$ & $0.638\pm0.038$ & Yes \\
\bottomrule
\end{tabular}
\par\smallskip
\textit{(b) Information}\par\smallskip
\begin{tabular}{@{}lccc@{}}
\toprule
Group & Raw information & Projected information & Surviving fraction \\
\midrule
Landbird/land & $0.737\pm0.046$ & $0.436\pm0.045$ & $0.590\pm0.027$ \\
Landbird/water & $0.794\pm0.032$ & $0.289\pm0.014$ & $0.365\pm0.026$ \\
Waterbird/land & $0.973\pm0.033$ & $0.764\pm0.031$ & $0.786\pm0.036$ \\
Waterbird/water & $2.063\pm0.054$ & $0.511\pm0.021$ & $0.248\pm0.015$ \\
\bottomrule
\end{tabular}
\par\smallskip
\end{table}


\paragraph{Foreground-preserving construction.}
Waterbirds combines CUB birds~\citep{wah2011cub} with Places backgrounds~\citep{zhou2018places} under soft CUB mattes. For each bird, we retain background pixels verbatim outside the dilated foreground and fill the covered region using Telea inpainting~\mbox{\citep{telea2004image}}, with a 5-pixel inpainting radius and a $5\times5$ elliptical dilation kernel. Its extracted bird layer is recomposed onto training-image donors of the same or opposite land/water background type. Donors have foreground fraction at most $0.08$; a deterministic sequence tries four, taking the first with visible inpainted area at most $0.03$, or the least-inpainted of four. The same-type arm undergoes identical recomposition.

Both arms use the same foreground layer, matte, label and source birds. Opaque foreground pixels are unchanged before JPEG encoding (quality 95, no chroma subsampling); soft-edge mixtures vary with background. Two draws per arm yield 23,176 composites. Visible inpainted area has mean $4.13\%$, p95 $8.45\%$, maximum $20.37\%$, and arm means $4.13\%$/$4.12\%$. Correctness is averaged across donors within bird, then across birds.

\paragraph{Outcomes and confident flips.}
Table~\ref{tab:d-intervention-outcomes} reports all groups. The confident half uses the within-slice median calibrated true-label probability among originally correct birds, including ties. In WW, opposite-type/control flip rates are $56.4\%$/$3.3\%$ overall and $41.8\%$/$0.7\%$ in that half; the fixed threshold $p\ge0.9$ selects no WW birds.

\begin{table}[htbp]
\centering\small
\setlength{\tabcolsep}{3pt}
\caption{Held-out interventions. Accuracies and flip rates are seed means in percent; controlled drops are same-type minus opposite-type percentage points with nominal 95\% normal-approximation half-widths. Panel (b) counts are mean confident-half counts.}
\label{tab:d-intervention-outcomes}
\label{tab:d-confident-flips}
\textit{(a) Paired intervention arms}\par\smallskip
\begin{tabular}{@{}lrcccc@{}}
\toprule
Group & Birds & Original & Same type & Opposite type & Controlled drop \\
\midrule
Landbird/land & 2,255 & 98.9 & 99.3 & 72.6 & $26.7\pm1.8$ \\
Landbird/water & 2,255 & 71.3 & 72.3 & 99.2 & $-26.9\pm2.0$ \\
Waterbird/land & 642 & 42.3 & 40.6 & 92.3 & $-51.7\pm2.0$ \\
\textbf{Waterbird/water} & 642 & \textbf{92.0} & \textbf{93.2} & \textbf{40.6} & $\mathbf{52.6\pm1.4}$ \\
\bottomrule
\end{tabular}
\par\smallskip
\textit{(b) Opposite-type flips among originally correct predictions}\par\smallskip
\begin{tabular}{@{}lccc@{}}
\toprule
Group & \shortstack{All originally\\correct} & \shortstack{Confident\\half} & \shortstack{Confident-half\\count} \\
\midrule
Landbird/land & 26.8 & 14.8 & 1,122 \\
Landbird/water & 0.2 & 0.1 & 805 \\
Waterbird/land & 1.2 & 0.3 & 136 \\
Waterbird/water & 56.4 & 41.8 & 299 \\
\bottomrule
\end{tabular}
\par\smallskip
\end{table}


\paragraph{Matched selection comparisons.}
Reapplying eligibility within each seed leaves LL and WW. Minimum validation accuracy and maximum uncalibrated NLL, evaluated separately on this identical set, agree with the surviving-fraction selector in all five seeds. All-group comparisons answer a different question; landbird/water determines the model-level certificate (Table~\ref{tab:d-selection-comparison}).
% Frozen-coordinate reanalysis; full means/matrices remain in the artifact.
\paragraph{Centering and predictive Fisher.}
The post hoc check fixes the checkpoints, temperatures, validation coordinates $v=(z,b)^\top$, original scales, eligible sets and intervention outcomes. No model/probe is refitted. With $\langle\cdot\rangle_g$ the divisor-$n_g$ slice mean, compare
\[
J_g=\langle vv^\top\rangle_g=\Sigma_g+\mu_g\mu_g^\top,\quad
\Sigma_g=\langle(v-\mu_g)(v-\mu_g)^\top\rangle_g,\quad \mu_g=\langle v\rangle_g.
\]
Centering does not rescale. The local model $\operatorname{logit}p_{\theta,\xi}(x)=f_w(x)/\widehat T+\theta z(x)+\xi b(x)$ at zero has score $(Y-p)v$ with conditional mean zero even when $\mu_g\ne0$. Its predictive Fisher and the common fraction are
\[
I_g=\langle p(1-p)vv^\top\rangle_g,\qquad
r_\rho(A)=1-\frac{A_{zb}^2}{A_{zz}(A_{bb}+\rho)},\quad\rho=10^{-3}.
\]
Perturbations follow calibration and profile only the declared background coordinate. The normalized-weight comparator is $K_g=I_g/\bar q_g$, $\bar q_g=\langle p(1-p)\rangle_g$, so $r_\rho(K_g)=r_{\bar q_g\rho}(I_g)$: normalization preserves the unregularized fraction but changes effective ridge. Table~\ref{tab:d-matrix-sensitivity} therefore reports it separately.
\begin{table}[htbp]
\centering\small
\setlength{\tabcolsep}{3pt}
\caption{Selection and matrix sensitivity, evaluated within each seed. Panel (a) distinguishes original eligible and all-group sets; drops are mean percentage points. Panel (b) fixes the original eligible LL/WW set and interventions; fractions/drops are means $\pm$ nominal 95\% normal-approximation half-widths.}
\label{tab:d-selection-comparison}
\label{tab:d-matrix-sensitivity}
\textit{(a) Matched selectors}\par\smallskip
\begin{tabular}{@{}p{.39\linewidth}p{.13\linewidth}p{.30\linewidth}r@{}}
\toprule
Statistic & Candidates & Selected group, seeds & Drop \\
\midrule
Minimum raw task information & Eligible & Landbird/land, 5/5 & 26.7 \\
Minimum absolute projected information & Eligible & Landbird/land, 5/5 & 26.7 \\
Minimum surviving fraction & Eligible & Waterbird/water, 5/5 & 52.6 \\
Minimum validation accuracy or maximum validation NLL & Eligible & Waterbird/water, 5/5 & 52.6 \\
Minimum normalized-weight fraction & Eligible & Waterbird/water, 5/5 & 52.6 \\
Maximum NLL or minimum accuracy & All four & Waterbird/land, 5/5 & $-51.7$ \\
Group determining $I_{0.95}^{\FG}$ & All four & Landbird/water, 5/5 & $-26.9$ \\
\bottomrule
\end{tabular}
\par\smallskip
\textit{(b) Frozen-coordinate constructions}\par\smallskip
\begin{tabular}{@{}lcccc@{}}
\toprule
Construction & LL fraction & WW fraction & Selected, seeds & Drop \\
\midrule
Uncentered $J_g$ & $0.590\pm0.027$ & $0.248\pm0.015$ & WW, 5/5 & $52.6\pm1.4$ \\
Centered $\Sigma_g$ & $0.779\pm0.028$ & $0.855\pm0.032$ & LL, 5/5 & $26.7\pm1.8$ \\
Predictive Fisher $I_g$ & $0.987\pm0.014$ & $0.273\pm0.016$ & WW, 5/5 & $52.6\pm1.4$ \\
Normalized weights $K_g$ & $0.987\pm0.014$ & $0.269\pm0.016$ & WW, 5/5 & $52.6\pm1.4$ \\
\bottomrule
\end{tabular}
\par\smallskip
\end{table}

Centering selects LL in all seeds, whereas predictive Fisher retains WW. Slice means therefore contribute to the original ranking: for WW in seed 22, $\mu_z^2=2.094$ versus $\Sigma_{zz}=0.0636$. Table~\ref{tab:d-matrix-per-seed} preserves the per-seed reversal; full means, covariance entries, weights, matrices and hashes are retained in the artifact.

\begin{table}[htbp]
\centering\small
\caption{Per-seed fractions on the original global scales. Each cell is LL / WW; the smaller value determines selection. Centering reverses the ranking in every seed.}
\label{tab:d-matrix-per-seed}
\begin{tabular}{@{}rccc@{}}
\toprule
Seed & Uncentered $r_\rho(J_g)$ & Centered $r_\rho(\Sigma_g)$ & Predictive $r_\rho(I_g)$ \\
\midrule
22 & 0.5570 / 0.2291 & 0.7430 / 0.8542 & 0.9817 / 0.2587 \\
23 & 0.6048 / 0.2494 & 0.8098 / 0.8481 & 0.9920 / 0.2729 \\
24 & 0.6356 / 0.2755 & 0.8132 / 0.8359 & 0.9620 / 0.3032 \\
25 & 0.5803 / 0.2392 & 0.7514 / 0.9156 & 0.9998 / 0.2571 \\
26 & 0.5716 / 0.2460 & 0.7779 / 0.8189 & 0.9998 / 0.2732 \\
\bottomrule
\end{tabular}
\end{table}


\paragraph{Raw, projected, and tail summaries.}
We compare four summaries of retained validation geometry: mean raw task second moment $A$, arithmetic mean projected information $P$, reciprocal mean inverse information $H$, and inverse-CVaR certificate $C$. For the four equally weighted groups, let $F_g=\widehat F_g^{\rid}$ and $L_g=(F_g+\varepsilon)^{-1}$. Then
\[
\begin{aligned}
A&=\tfrac14\sum_g\widehat J_{TT,g}, &
P&=\tfrac14\sum_g F_g,\\
H&=\left(\tfrac14\sum_g L_g\right)^{-1}, &
C&=\left[\CVaR_{0.95,\pi}(L_g)\right]^{-1},\qquad \pi_g=\tfrac14.
\end{aligned}
\]
Thus $H$ is stabilized harmonic projected information, not mean inverse information. All four summaries use uncentered feature moments, nuisance ridge $10^{-3}$, and reciprocal stabilization $10^{-8}$. This post hoc analysis joins validation summaries to held-out interventions in the primary (seeds 22--26) and development (17--21) cohorts. It compares summaries of existing models, not training objectives.

\begin{table}[htbp]
\centering\small
\setlength{\tabcolsep}{3pt}
\caption{Four validation summaries and their intervention associations. Panels (a)--(b): means $\pm$ nominal 95\% normal-approximation half-widths, $1.96s/\sqrt5$; changes are paired within seed and drops are signed percentage points. Selected FG includes accuracy-budget fallbacks (primary seeds 22, 24; development seed 20). Panel (c): descriptive Spearman coefficients for ten models arranged in five ERM/FG seed pairs per cohort, with average ranks for tied outcomes.}
\label{tab:four-way-swap}
\textit{(a) Model summaries and paired changes}\par\smallskip
\begin{tabular}{@{}llcccc@{}}
\toprule
Cohort & Model/change & $A$: raw & $P$: projected & $H$: harmonic & $C$: inverse-CVaR\\
\midrule
Primary & ERM & $1.142\pm0.012$ & $0.500\pm0.010$ & $0.442\pm0.011$ & $0.289\pm0.014$\\
& Selected FG & $0.730\pm0.086$ & $0.484\pm0.047$ & $0.435\pm0.043$ & $0.330\pm0.043$\\
& FG $-$ ERM & $-0.412\pm0.079$ & $-0.016\pm0.046$ & $-0.007\pm0.041$ & $+0.041\pm0.046$\\
Development & ERM & $1.152\pm0.005$ & $0.489\pm0.011$ & $0.413\pm0.014$ & $0.255\pm0.018$\\
& Selected FG & $0.754\pm0.118$ & $0.519\pm0.065$ & $0.433\pm0.057$ & $0.303\pm0.082$\\
& FG $-$ ERM & $-0.398\pm0.118$ & $+0.030\pm0.063$ & $+0.020\pm0.066$ & $+0.048\pm0.095$\\
\bottomrule
\end{tabular}
\par\smallskip
\textit{(b) ERM group-selection rules and controlled drops}\par\smallskip
\begin{tabular}{@{}>{\raggedright\arraybackslash}p{.34\linewidth}llrr@{}}
\toprule
Rule & Candidates & Selected, each cohort & Primary & Development\\
\midrule
Minimum raw second moment & All four & LL 4/5; LW 1/5 & $+15.8\pm21.3$ & $+16.9\pm23.2$\\
Minimum projected / largest inverse contribution / CVaR-active group & All four & LW 5/5 & $-26.9\pm2.0$ & $-29.7\pm2.2$\\
Largest inverse contribution & Eligible & LL 5/5 & $+26.7\pm1.8$ & $+29.2\pm2.8$\\
\bottomrule
\end{tabular}
\par\smallskip
\textit{(c) Model-level rank associations with held-out outcomes}\par\smallskip
\begin{tabular}{@{}lrrrr@{}}
\toprule
& \multicolumn{2}{c}{Primary} & \multicolumn{2}{c}{Development}\\
Summary & $D_{\max}$ & $W_{\mathrm{opp}}$ & $D_{\max}$ & $W_{\mathrm{opp}}$\\
\midrule
$A$: mean raw & $-0.10$ & $+0.26$ & $+0.26$ & $-0.13$\\
$P$: mean projected & $-0.09$ & $+0.21$ & $-0.03$ & $-0.43$\\
$H$: reciprocal mean inverse & $-0.15$ & $+0.21$ & $-0.12$ & $-0.40$\\
$C$: inverse-CVaR & $+0.15$ & $-0.24$ & $-0.03$ & $-0.48$\\
\bottomrule
\end{tabular}
\end{table}

For Table~\ref{tab:four-way-swap}(c), $D_{\max}=100\max_g(\operatorname{Acc}_{\mathrm{same},g}-\operatorname{Acc}_{\mathrm{opposite},g})$ and $W_{\mathrm{opp}}=\min_g\operatorname{Acc}_{\mathrm{opposite},g}$, both over all four groups. Integer correct counts recover exact outcome ties: primary ERM seeds 22 and 25 share drop $348/642$, despite different floating-point subtractions. LL and LW denote landbird/land and landbird/water. Negative drops mean opposite-type recomposition improves accuracy. The largest group contributes $38.3\%$ and $40.6\%$ of summed inverse information on average in the primary and development cohorts, respectively. The original headline rules differed: surviving fraction in the primary cohort, absolute projected information in development; this comparison consistently uses all four groups for model summaries and correlations.

Projection and aggregation give different model summaries, with rank associations varying across cohorts. For four equally weighted groups and $\alpha=0.95$, $C=\min_gF_g+\varepsilon$; its active group also minimizes projected information and maximizes inverse contribution. Here that group is LW, whose opposite-type recomposition improves accuracy. The prospective WW diagnosis instead uses the surviving fraction within the eligible set. This comparison does not establish superior swap-damage prediction by inverse-CVaR; paired changes in $P,H,C$ all have intervals containing zero. The artifact provides per-model and per-seed values, exact-count outcomes, input hashes, and replay code.


\paragraph{Auxiliary flip prediction.}
Originally correct predictions contribute two opposite-type donor rows per bird. Baseline predictors are uncalibrated NLL, calibrated true-label confidence, absolute raw logit and standardized task logit $z$. Augmentation adds $\beta_gb$ and $|z-\beta_gb|$, with $\beta_g=\widehat J_{TN,g}/(\widehat J_{NN,g}+10^{-3})$ frozen from the validation block. Both models use identical five-fold, nonshuffled CUB-species \texttt{GroupKFold} splits, keeping species and donor pairs together. Logistic regression uses an intercept, $\ell_2$ penalty, $C=1$, L-BFGS, tolerance $10^{-4}$, maximum 5,000 iterations, seed 0 and no class weighting or tuning. All columns, including group indicators, are centered and scaled using training-fold means and population standard deviations; standard deviations below $10^{-9}$ are replaced by one.

Pool out-of-fold predictions within seed before computing average precision (AUPRC), AUROC and mean log likelihood. For clipped probabilities $q\in[10^{-12},1-10^{-12}]$,
\begin{equation}
G_s=\frac1{N_s}\sum_{i=1}^{N_s}\!\left[\operatorname{LL}(y_{si},q^A_{si})-\operatorname{LL}(y_{si},q^B_{si})\right],\quad
\operatorname{LL}(y,q)=y\log q+(1-y)\log(1-q).
\label{eq:d-flip-loglik-gain}
\end{equation}
$G_s$ is augmented-minus-baseline gain in nats per donor row. Table~\ref{tab:d-flip-prediction} averages paired changes equally across seeds 22--26. The coordinates improve prediction without group indicators and within WW; with indicators, AUPRC decreases. This tests their incremental value, not unique attribution to the scalar information statistic. Original inputs, exact-source refits and paired records are in the artifact.

\begin{table}[htbp]
\centering\small
\setlength{\tabcolsep}{3pt}
\caption{Incremental flip prediction. Counts and rates are seed means; parentheses count seeds with improvement. Changes are paired within seed.}
\label{tab:d-flip-prediction}
\begin{tabular}{@{}p{.24\linewidth}rrccc@{}}
\toprule
Analysis & Rows & \shortstack{Flip\\rate} & \shortstack{Log-likelihood\\gain} & $\Delta$AUPRC & $\Delta$AUROC \\
\midrule
All groups, no indicators & 9,398 & 0.199 & $+0.016$ & $+0.035$ (5/5) & $+0.096$ (5/5) \\
All groups, with indicators & 9,398 & 0.199 & $+0.001$ & $-0.007$ (0/5) & $+0.003$ (5/5) \\
Selected slice & 1,181 & 0.564 & $+0.013$ & $+0.026$ (5/5) & $+0.015$ (5/5) \\
\bottomrule
\end{tabular}
\par\smallskip
\end{table}


\paragraph{Development chronology.}
The surviving-fraction rule followed an absolute-information pilot on seeds 17--21 and was frozen for newly trained seeds 22--26 on the same benchmark. These tables consistently use the documented primary rerun; the original-versus-rerun ledger is retained in the artifact. Appendix~\ref{app:extended-training} gives the Fisher-Glass training comparison under this same intervention.

\FloatBarrier
\subsection{Synthetic Classification and Tail Checks}
\label{app:synthetic-checks}

Draw $U_e\sim\mathrm{Uniform}(0,1)$, set $\delta_e=0.35U_e^{1/\gamma}$, and draw independent standard normals $s,r$ for each observation. The task is $z=s+\delta_er$, $Y=\mathbf1\{z>0\}$, scored by nuisance-only logit $4s$. Thus
\begin{equation}
\Pr(\delta_e<t)=(t/0.35)^\gamma\ (0\le t\le0.35),\qquad J_{TT,e}=1+\delta_e^2,\quad F_e^0=\delta_e^2.
\label{eq:d-synthetic-generator}
\end{equation}
Population feature moments are exact; loss/accuracy use 256 observations per environment. The $\gamma=1$ experiment uses 300,000 environments and seeds 11--15. Fit the log-survival slope over the largest 5\% of inverse information; empirical CVaR averages values at or above the empirical $\alpha$-quantile. Figure~\ref{fig:geometry-tails} regenerates seed 11 with $L_e=(\delta_e^2+10^{-12})^{-1}$, rather than recovering missing original raw output; its survival curve and the five-seed summaries have different replication levels.

Table~\ref{tab:d-synthetic-summary} checks Theorem~\ref{thm:inverse-information-geometry} and the scalar portfolio law in Appendix~\ref{app:theory-tails}. For $K$ independent draws, information $K^{-1}\sum_j\delta_{e,j}^2$ has exponent $K/2$ at $\gamma=1$; weights affect its constant. The mass sweep predicts $\gamma/2$. An anchor with $\delta_0=0.5$ and reserve mass $0.05$ yields information $0.95\delta_e^2+0.05(0.25)$, floor $0.0125$, and measured certificate $0.01260\pm0.000001$ (five seeds; Appendix~\ref{app:information-reserves}).

\begin{table}[htbp]
\centering\small
\setlength{\tabcolsep}{3pt}
\caption{Synthetic checks. Panel (a) and the branch rows are five-seed means $\pm$ nominal 95\% normal-approximation half-widths; zero widths reflect rounding. Mass-sweep rows are reported single-seed values.}
\label{tab:d-synthetic-summary}
\label{tab:d-structural-comparisons}
\textit{(a) Prediction and information}\par\smallskip
\begin{tabular}{@{}lc@{}}
\toprule
Quantity & Estimate \\
\midrule
Mean accuracy & $94.54\pm0.00\%$ \\
Fifth-percentile environment accuracy & $88.67\pm0.00\%$ \\
$\CVaR_{0.95}$ of environment-mean NLL & $0.262\pm0.000$ \\
$I_{0.95}^{\FG}$ & $(1.22\pm0.53)\times10^{-8}$ \\
Inverse-information exponent $\widehat\beta$ & $0.498\pm0.003$ \\
\bottomrule
\end{tabular}
\par\smallskip
\textit{(b) Inverse-information tail exponents}\par\smallskip
\begin{tabular}{@{}lccc@{}}
\toprule
Construction & Setting & Predicted $\beta$ & Estimated $\beta$ \\
\midrule
Mass near the null & $\gamma=1$ & 0.5 & 0.509 \\
Mass near the null & $\gamma=2$ & 1 & 0.985 \\
Mass near the null & $\gamma=4$ & 2 & 2.001 \\
Independent branches & $K=1$ & 0.5 & $0.499\pm0.003$ \\
Independent branches & $K=2$ & 1 & $1.000\pm0.009$ \\
Independent branches & $K=3$ & 1.5 & $1.508\pm0.015$ \\
\bottomrule
\end{tabular}
\par\smallskip
\end{table}


\FloatBarrier
\subsection{Hidden-Slice Discovery Results}
\label{app:discovery-results}

Appendix~\ref{app:hidden-slice-discovery} gives the blind protocol: task labels, frozen MobileNet features and a fixed CLIP direction are available; background/group metadata remain sealed through held-out evaluation. Results use the primary hidden-slice cohort, seeds 17--21, with full-validation gates frozen before test. Support correction may depend on the test sample but leaves all five primary CLIP gates' test weights unchanged. Information uses same-class pooling; accuracy, uncalibrated NLL and composition use admissible within-class weights without pooling. The artifact identifies this cohort as v2 and records original hashes; the older uncorrected cohort is excluded.

All selectors share classifiers with average accuracy $76.3\pm0.7\%$ and post hoc worst-group accuracy $57.3\pm1.3\%$. CLIP search finds $98.8\%$ membership-weighted slice accuracy. Table~\ref{tab:d-hidden-selectors} includes every matched-mass comparator and paired contrast. The random-direction interval and the uncentered and centered feature-moment intervals against high-loss selection contain zero. Feature controls reevaluate gates selected using predictive Fisher without rediscovery.

\begin{table}[htbp]
\centering\small
\setlength{\tabcolsep}{3pt}
\caption{Hidden-slice comparisons. Panel (a): means $\pm$ 95\% Student-$t$ half-widths across five paired seeds; predictive information retains calibration and shrinkage $\eta=0.10$. Panel (b): CLIP-search minus comparator, with paired 95\% Student-$t$ intervals; negative differences indicate lower information.}
\label{tab:d-hidden-selectors}
\label{tab:d-hidden-contrasts}
\textit{(a) Held-out selectors}\par\smallskip
\begin{tabular}{@{}p{.53\linewidth}cc@{}}
\toprule
Selector & \shortstack{Predictive\\projected information} & \shortstack{Uncalibrated\\slice NLL} \\
\midrule
Fisher search, fixed CLIP direction & $0.005\pm0.001$ & $0.031\pm0.017$ \\
High loss, same target within-class mass & $0.019\pm0.006$ & $5.779\pm0.496$ \\
Random, same target within-class mass & $0.013\pm0.001$ & $1.025\pm0.059$ \\
Constant class-conditional gate & $0.013\pm0.001$ & $1.027\pm0.049$ \\
Fisher search, random direction & $0.006\pm0.001$ & $0.031\pm0.015$ \\
\bottomrule
\end{tabular}
\par\smallskip
\textit{(b) Paired information contrasts}\par\smallskip
\begin{tabular}{@{}llcc@{}}
\toprule
Information & Comparison & Difference & 95\% interval \\
\midrule
Predictive & High-loss slice & $-0.014$ & $[-0.020,-0.008]$ \\
Predictive & Random slice & $-0.008$ & $[-0.009,-0.007]$ \\
Predictive & Constant gate & $-0.008$ & $[-0.009,-0.007]$ \\
Predictive & Random direction & $-0.0004$ & $[-0.0015,0.0007]$ \\
Unweighted feature & High-loss slice & $-0.034$ & $[-0.106,0.039]$ \\
Centered unweighted & High-loss slice & $-0.071$ & $[-0.166,0.024]$ \\
\bottomrule
\end{tabular}
\par\smallskip
\end{table}


\paragraph{Class selection and overlap.}
CLIP, random-direction, random-slice and constant-gate selectors choose landbirds in all five seeds; high loss chooses landbirds in seeds 17,18,21 and waterbirds in 19,20. The two different-class Jaccards are zero; same-class values are $0.00222$, $0.00111$, and $0$. Their mean $0.000666$, reported as $0.001$, combines between-class and within-class differences under Appendix~\ref{app:discovery-crossfit}'s top-count rule.

\paragraph{Support and sensitivity.}
On same-class rows from the 15\% training-calibration subset, cross-fold gate-score correlation is $0.419\,[0.353,0.485]$ and overlap $0.258\,[0.225,0.291]$. Test ESS is $414\,[383,445]$ using admissible weights. Raw membership mass is $0.083\,[0.077,0.089]$ within the selected class (target $0.10$), not the full test set; landbirds contain 4,510 test examples. Validation cross-fitting, separate from final testing, gives information $0.004$--$0.005$ over masses $\{0.05,0.10,0.20\}$ versus high-loss $0.012$--$0.027$, and $0.004$--$0.006$ over shrinkages $\{0.05,0.10,0.20\}$.

Composition after held-out evaluation is $97.0\%$ landbird/land for CLIP, $94.4\%$ for random direction, and $58.4\%$ landbird/water for high loss. The hidden slice has lower projected predictive information than the weakest predefined group in all five seeds: mean paired difference $-0.0029$, with 95\% Student-$t$ interval $[-0.0034,-0.0024]$, taking the predefined-group minimum separately within each seed. Both sides use the same model, CLIP coordinate, calibrated predictive moments, pooling $\eta=0.10$, and ridge $10^{-3}$. A metadata-trained oracle probe used only after blind freeze gives information $0.003\pm0.000$, NLL $0.009\pm0.004$, and $99.1\%$ landbird/land composition (zero width is rounding). Repair outcomes are in Appendix~\ref{app:extended-training}.

\FloatBarrier
\subsection{Predefined-Group Diagnostics}
\label{app:predefined-diagnostics}

\paragraph{Fitting.}
CivilComments-WILDS~\citep{borkan2019nuanced,koh2021wilds} predicts comment toxicity, with $y=0$ denoting nontoxic comments. Identity-any indicates that at least one of the 24 original identity annotation scores is at least $0.5$. Training-fitted unigram/bigram TF-IDF (50,000 features, minimum document frequency 5, sublinear TF) feeds class-balanced SGD logistic task/identity-any probes ($\ell_2=10^{-5}$, 50 passes, tolerance $10^{-3}$). Identity-by-label groups require 100 examples; tail regression uses the largest 50\% of group inverse information, requiring eight groups and four tail points. Seed 13 gives the fixed diagnostic; seeds 13--17 give aggregates.

Waterbirds uses frozen ImageNet MobileNetV3-small features (extraction batch 64) and balanced L-BFGS logistic task/background models (1,000 maximum iterations). Four groups are reported directly, without fitting a tail exponent; only controls/anchors vary over seeds 17--21. Primary task/probe and metadata coordinates use training moments. These original runs used an i9-11900F/32 GB workstation, about one minute for CivilComments and three minutes for first Waterbirds extraction/fitting, then cached features.

\begin{table}[htbp]
\centering\small
\setlength{\tabcolsep}{3pt}
\caption{Original predefined-group diagnostics, distinct from fresh audits in Appendix~\ref{app:implemented-audit}. Means $\pm$ nominal 95\% normal-approximation half-widths; Waterbirds zeros reflect fixed-model fitting.}
\label{tab:d-predefined-diagnostics}
\begin{tabular}{@{}lcccc@{}}
\toprule
Dataset & \shortstack{Average\\accuracy (\%)} & \shortstack{Worst-group\\accuracy (\%)} & \shortstack{Loss-\\$\CVaR_{0.95}$} & $I_{0.95}^{\FG}$ \\
\midrule
CivilComments & $86.10\pm0.22$ & $43.75\pm0.64$ & $0.998\pm0.012$ & $0.500\pm0.001$ \\
Waterbirds & $77.11\pm0.00$ & $56.54\pm0.00$ & $1.741\pm0.000$ & $0.165\pm0.000$ \\
\bottomrule
\end{tabular}
\par\smallskip
\end{table}


For fixed CivilComments, the weakest-information group is \texttt{muslim}/$y=0$ (accuracy $60.1\%$); lowest accuracy is \texttt{black}/$y=0$ ($44.0\%$). Waterbirds' weakest-information group is landbird/land ($97.3\%$ accuracy), whereas waterbird/land has lowest accuracy ($56.5\%$); inverse-information p95 is $5.90$.

\paragraph{Specification and stress controls.}
The prespecified semantic family is $\{b_{\rm meta},b_{\rm probe}\}$: standardized identity-any/background annotation and its training-fitted probe, without cross-fitting. Semantic relevance follows those annotations; no held-out relevance threshold, anchor task-preservation tolerance, fold-stability admission screen, or nuisance-conditioning threshold was applied. The ridge is $10^{-3}$; family minima are $0.499$ and $0.165$. Reserve anchors (identity-present label-balanced CivilComments; four-group-balanced Waterbirds) are construction inputs, not admission samples. Fresh exact/ridge audits and hidden-gate stability are separate analyses.

\begin{table}[htbp]
\centering\small
\setlength{\tabcolsep}{3pt}
\caption{Nuisance sensitivity of fixed-model diagnostics. Entries are $I_{0.95}^{\FG}$; worst-group accuracy remains $44.0\%$/$56.5\%$ for CivilComments/Waterbirds.}
\label{tab:d-nuisance-sensitivity}
\begin{tabular}{@{}lccp{.27\linewidth}@{}}
\toprule
Specification & CivilComments & Waterbirds & Role \\
\midrule
Metadata & 0.587 & 0.266 & Semantic-family member \\
Learned probe & 0.499 & 0.165 & Semantic-family minimum \\
Random & 0.588 & 0.482 & Negative control \\
Under-projected & 0.588 & 0.482 & Under-projection control \\
Over-projected & 0.001 & 0.012 & Task-erasure control \\
Hidden-loss & 0.060 & 0.177 & Loss-coordinate stress control \\
\bottomrule
\end{tabular}
\par\smallskip
\end{table}


Random nuisance is independent standard normal; under-projection sets it to zero. Over-projection uses standardized task score plus Gaussian noise of standard deviation $0.01$ and deliberately erases task information. ``Hidden-loss'' is observed test loss with original groups fixed. The random, over-projected, and observed-loss coordinates use evaluated-array normalization and are not independently learned nuisance scores. The metadata/random contrast is $0.587$/$0.588$ on CivilComments, versus probe/random $0.165$/$0.482$ on Waterbirds, illustrating specification dependence.

The versioned artifact (see the \hyperref[sec:artifact-availability]{artifact availability statement}) retains execution commands, full-precision matrices and overlaps, run-history ledgers and expanded source tables, with original cohort provenance. Those records supplement the scientific definitions and controls above.
\FloatBarrier
